\section{Conclusions}
Agentic recommender systems represent a significant shift from passive item-ranking pipelines to autonomous, goal-directed systems capable of reasoning, planning, acting, and adapting within complex environments. 
By integrating large language model agents with profiles, memory structures, tool interfaces, and multi-step workflows, modern systems move beyond traditional assumptions about static user intents and fixed recommendation stages. Our survey organizes this emerging field through three complementary paradigms, \ie agent-assisted recommendation, agent-as-recommender, and agent-as-user-simulator. This survey situates them within a unified autonomy framework that spans retrieval augmentation, tool use, single-agent planning, and multi-agent collaboration.

Across studies, we observe clear trends. Agent-assisted systems demonstrate strong gains in semantic understanding, retrieval quality, and contextual relevance. Single-agent recommenders further introduce coherent reasoning, long-horizon planning, and dynamic preference modeling, while multi-agent frameworks enable division of labor, critique, negotiation, and consensus building. User simulation agents provide new opportunities for scalable evaluation, counterfactual analysis, and environment modeling, allowing the field to explore settings that are difficult to observe in real-world logs. 
Despite these advances, substantial challenges remain. Lifelong user modeling, contextual abstraction, and multimodal grounding are still far from solved. User experience issues such as controllability, explainability, trustworthiness, and privacy require dedicated attention as systems become more autonomous. At the system level, scalability, efficiency, and robust agent–tool orchestration will determine whether agentic approaches can be deployed at industrial scale. Addressing these challenges creates opportunities to propose new tasks, design specialized modules, and introduce capabilities that reshape the user experience of recommendation.